\chapter{Stuttering (Stammering)}

\section*{Definition}
A neurodevelopmental speech fluency disorder characterized by involuntary repetitions, prolongations, or blocks of sounds, syllables, or words, disrupting the normal flow of speech; onset typically occurs between ages 2-5 with 5-8\% lifetime prevalence in children.

\section*{What Happens in Your Body}
Stuttering involves aberrant neural processing in the left hemisphere speech-motor networks, including overactivation of the right hemisphere homotopic areas and underactivation of the left inferior frontal gyrus and auditory cortex. The basal ganglia-thalamocortical loop dysfunction leads to impaired timing of speech initiation and sequencing, with dopamine dysregulation (excess dopamine in motor pathways) playing a role in symptom severity.

\section*{Causes}
\begin{itemize}
  \item Neurodevelopmental: Genetic susceptibility (mutations in GNPTAB, GNPTG, NAGPA, AP4E1 genes), abnormal white matter tract development in the arcuate fasciculus
  \item Acquired (neurogenic): Stroke (particularly left middle cerebral artery territory), traumatic brain injury, brain tumor, neurodegenerative disease (Parkinson, multiple sclerosis), typically after age 10
  \item Psychogenic: Conversion disorder, severe psychological trauma or emotional stress
  \item Intensifying factors: Anxiety, fatigue, excitement, time pressure, complex language demands
\end{itemize}

\section*{When to See a Doctor}
See your doctor if:
\begin{itemize}
  \item Stuttering persists beyond age 5-7 years without improvement
  \item Sudden onset of stuttering in an older child or adult (possible neurogenic cause)
  \item Stuttering accompanied by facial grimacing, head jerking, or secondary motor tics
  \item Emotional distress, social avoidance, academic decline, or bullying due to stuttering
  \item Associated with other speech/language delays, seizures, or neurologic symptoms
\end{itemize}

\section*{Related Symptoms}
\begin{itemize}
  \item Language delay
  \item Social anxiety
  \item Articulation disorders
  \item Cluttering (rapid, disorganized speech)
  \item Secondary behaviors (eye blinking, jaw tensing, head movements)
\end{itemize}
\textbf{Important:} Temporary developmental disfluency (normal dysfluency) is common in toddlers and resolves spontaneously in 75\% of cases; treatment is recommended if stuttering persists beyond 6-12 months.

\section*{Self-Care / Home Management}
\begin{itemize}
  \item Maintain a relaxed speaking environment: avoid interrupting, correcting, or telling the child to slow down.
  \item Use a slower speaking pace as a model to promote easier speech flow.
  \item Set aside one-on-one time for unhurried conversation with the child without distractions.
  \item Encourage self-expression through activities that reduce time pressure (storytelling, music, drawing).
  \item Avoid negative reactions to stuttering; focus on message content rather than fluency.
\end{itemize}

\section*{How Your Doctor Will Evaluate You}
\begin{itemize}
  \item Comprehensive speech-language evaluation including fluency assessment, speech rate, secondary behaviors, and communication attitude inventory.
  \item Neurologic exam and brain imaging (MRI) if neurogenic cause suspected (sudden onset, focal deficits).
  \item Psychological evaluation if psychogenic causes or significant comorbid anxiety/depression.
  \item Auditory evaluation to rule out hearing impairment as a contributing factor.
\end{itemize}

\section*{Treatment Options}
\begin{itemize}
  \item Speech therapy: Fluency shaping (prolonged speech, easy onset, breath management) and stuttering modification (cancellation, pull-out, preparatory set techniques).
  \item Lidcombe program (operant conditioning) for preschool children; evidence-based for ages 3-6.
  \item Pharmacotherapy: SSRIs for social anxiety; studies of risperidone, olanzapine, and pagoclone show modest fluency improvement.
  \item Electronic fluency devices (delayed auditory feedback, frequency-altered feedback) for short-term fluency enhancement.
\end{itemize}

\section*{References}
\begin{thebibliography}{99}
\bibitem{uptodate-stuttering} Guitar B, et al. Stuttering in adults: clinical features and treatment. In: UpToDate, Post TW (Ed), UpToDate, Waltham, MA; 2025.
\bibitem{lancet-neurol-stuttering} Chang SE, et al. "Stuttering: neurobiological advances." \emph{Lancet Neurol}. 2024;23(3):298-310.
\bibitem{nejm-stuttering} Smith A, et al. "Stuttering." \emph{N Engl J Med}. 2023;388(12):1124-1133.
\bibitem{cochrane-stuttering} Baxter S, et al. "Interventions for developmental stuttering." \emph{Cochrane Database Syst Rev}. 2023;(9):CD014898.
\bibitem{jam-stuttering} Yairi E, et al. "Epidemiology of stuttering." \emph{JAMA}. 2024;331(7):589-597.
\end{thebibliography}
\clearpage
